% ============================================================
\chapter{Containerization --- Docker for ML}
\label{ch:docker}
\index{Docker}
% ============================================================

\section{Why Containers?}
\label{sec:why-containers}

Virtual environments solve Python dependency isolation, but they do
\emph{not} capture system-level dependencies: OS libraries, CUDA
drivers, compilers, or specific tool versions. Containers solve this
by packaging an entire filesystem snapshot that runs identically
everywhere.

\begin{definition}[Container]
A \textbf{container}\index{container} is a lightweight, standalone,
executable unit of software that includes everything needed to run an
application: code, runtime, system tools, libraries, and settings.
Unlike virtual machines, containers share the host kernel, making them
much faster to start and smaller in footprint.
\end{definition}

\begin{center}
\begin{tikzpicture}[
  box/.style={rectangle, draw, rounded corners, align=center,
    font=\small, minimum width=2.8cm},
  layer/.style={rectangle, draw, minimum height=0.7cm, align=center,
    font=\scriptsize},
  arr/.style={-{Stealth[length=2mm]}, thick}
]
  % VM side
  \node[font=\small\bfseries] (vmtitle) {Virtual Machines};
  \node[layer, fill=gray!20, minimum width=6cm, below=0.2cm of vmtitle]
    (hw1) {Infrastructure (Hardware)};
  \node[layer, fill=orange!20, minimum width=6cm, above=0cm of hw1.north]
    (hyp) {Hypervisor};
  \node[layer, fill=blue!10, minimum width=1.7cm, above=0cm of hyp.north,
    xshift=-2cm] (os1) {Guest OS};
  \node[layer, fill=blue!10, minimum width=1.7cm, above=0cm of hyp.north]
    (os2) {Guest OS};
  \node[layer, fill=blue!10, minimum width=1.7cm, above=0cm of hyp.north,
    xshift=2cm] (os3) {Guest OS};
  \node[layer, fill=green!15, minimum width=1.7cm, above=0cm of os1.north]
    (a1) {App A};
  \node[layer, fill=green!15, minimum width=1.7cm, above=0cm of os2.north]
    (a2) {App B};
  \node[layer, fill=green!15, minimum width=1.7cm, above=0cm of os3.north]
    (a3) {App C};

  % Container side
  \node[font=\small\bfseries, right=3.5cm of vmtitle] (ctitle) {Containers};
  \node[layer, fill=gray!20, minimum width=6cm, below=0.2cm of ctitle]
    (hw2) {Infrastructure (Hardware)};
  \node[layer, fill=orange!10, minimum width=6cm, above=0cm of hw2.north]
    (hos) {Host OS};
  \node[layer, fill=cyan!15, minimum width=6cm, above=0cm of hos.north]
    (dock) {Container Engine (Docker)};
  \node[layer, fill=green!15, minimum width=1.7cm, above=0cm of dock.north,
    xshift=-2cm] (c1) {App A};
  \node[layer, fill=green!15, minimum width=1.7cm, above=0cm of dock.north]
    (c2) {App B};
  \node[layer, fill=green!15, minimum width=1.7cm, above=0cm of dock.north,
    xshift=2cm] (c3) {App C};
\end{tikzpicture}
\end{center}

\begin{remark}
Containers are \emph{not} lightweight VMs. They share the host kernel
and use Linux namespaces and cgroups for isolation. This makes them
start in milliseconds rather than minutes, and use megabytes rather
than gigabytes of overhead.
\end{remark}

\section{Docker Core Concepts}
\label{sec:docker-concepts}
\index{Docker!concepts}

\begin{keyformulas}[title=\textbf{Docker vocabulary}]
\begin{description}
  \item[Image] A read-only template with instructions for creating a
    container. Built from a \file{Dockerfile}.
  \item[Container] A running instance of an image. Ephemeral by default.
  \item[Layer] Each instruction in a Dockerfile creates a layer. Layers
    are cached and shared between images.
  \item[Registry] A repository for storing and distributing images
    (Docker Hub, GitHub Container Registry, AWS ECR).
  \item[Volume] Persistent storage that survives container destruction.
  \item[Dockerfile] A text file with instructions to assemble an image.
\end{description}
\end{keyformulas}

\begin{center}
\begin{tikzpicture}[
  box/.style={rectangle, draw, rounded corners, minimum width=2.5cm,
    minimum height=1cm, align=center, font=\small},
  arr/.style={-{Stealth[length=2.5mm]}, thick, steelblue}
]
  \node[box, fill=blue!10] (df) {\textbf{Dockerfile}\\(instructions)};
  \node[box, fill=green!10, right=2.5cm of df] (img) {\textbf{Image}\\(read-only)};
  \node[box, fill=orange!10, right=2.5cm of img] (ctr) {\textbf{Container}\\(running)};
  \node[box, fill=purple!10, above=1.2cm of img] (reg) {\textbf{Registry}\\(Docker Hub)};

  \draw[arr] (df) -- node[above, font=\scriptsize] {\cli{docker build}} (img);
  \draw[arr] (img) -- node[above, font=\scriptsize] {\cli{docker run}} (ctr);
  \draw[arr] (img) -- node[left, font=\scriptsize] {\cli{docker push}} (reg);
  \draw[arr] (reg) -- node[right, font=\scriptsize] {\cli{docker pull}} (img);
\end{tikzpicture}
\end{center}

\section{Essential Docker Commands}
\label{sec:docker-commands}

\begin{codeblock}[title=\textbf{Docker command reference}]
\begin{minted}{bash}
# Build an image from a Dockerfile
docker build -t my-ml-app:v1 .

# Run a container
docker run -p 8000:8000 my-ml-app:v1

# Run interactively with GPU support
docker run -it --gpus all my-ml-app:v1 bash

# List running containers
docker ps

# Stop and remove a container
docker stop <container_id>
docker rm <container_id>

# List images
docker images

# Remove an image
docker rmi my-ml-app:v1

# View container logs
docker logs -f <container_id>

# Execute a command in a running container
docker exec -it <container_id> bash
\end{minted}
\end{codeblock}

\section{Dockerfile for ML Projects}
\label{sec:dockerfile-ml}
\index{Dockerfile}

\begin{codeblock}[title=\textbf{Complete ML Dockerfile with Python and CUDA}]
\begin{minted}{dockerfile}
# Use NVIDIA CUDA base image for GPU support
FROM nvidia/cuda:12.1.0-cudnn8-runtime-ubuntu22.04

# Prevent interactive prompts during build
ENV DEBIAN_FRONTEND=noninteractive
ENV PYTHONUNBUFFERED=1
ENV PYTHONDONTWRITEBYTECODE=1

# Install system dependencies
RUN apt-get update && apt-get install -y --no-install-recommends \
    python3.11 \
    python3.11-venv \
    python3-pip \
    git \
    curl \
    && rm -rf /var/lib/apt/lists/*

# Set Python 3.11 as default
RUN ln -sf /usr/bin/python3.11 /usr/bin/python

# Create non-root user
RUN useradd -m -s /bin/bash mluser
WORKDIR /app

# Copy and install dependencies first (layer caching)
COPY requirements.txt .
RUN pip install --no-cache-dir -r requirements.txt

# Copy application code
COPY src/ ./src/
COPY models/ ./models/
COPY configs/ ./configs/

# Switch to non-root user
USER mluser

# Expose API port
EXPOSE 8000

# Health check
HEALTHCHECK --interval=30s --timeout=10s --retries=3 \
    CMD curl -f http://localhost:8000/health || exit 1

# Default command
CMD ["python", "-m", "uvicorn", "src.api:app", \
     "--host", "0.0.0.0", "--port", "8000"]
\end{minted}
\end{codeblock}

\begin{remark}
The order of \texttt{COPY} instructions matters for layer caching.
Dependencies change less often than source code, so
\file{requirements.txt} is copied and installed \emph{before} the
application code. This way, rebuilding after a code change does not
re-install all packages.
\end{remark}

\section{Multi-stage Builds}
\label{sec:multi-stage}
\index{Docker!multi-stage}

Multi-stage builds produce smaller production images by separating the
build environment from the runtime environment.

\begin{codeblock}[title=\textbf{Multi-stage Dockerfile for ML}]
\begin{minted}{dockerfile}
# ---------- Stage 1: Builder ----------
FROM python:3.11-slim AS builder

WORKDIR /build
COPY requirements.txt .
RUN pip install --no-cache-dir --prefix=/install \
    -r requirements.txt

# ---------- Stage 2: Training ----------
FROM python:3.11-slim AS trainer

WORKDIR /app
COPY --from=builder /install /usr/local
COPY src/ ./src/
COPY configs/ ./configs/
COPY data/ ./data/

RUN python src/train.py --config configs/config.yaml

# ---------- Stage 3: Production ----------
FROM python:3.11-slim AS production

RUN useradd -m mluser
WORKDIR /app

COPY --from=builder /install /usr/local
COPY --from=trainer /app/models/ ./models/
COPY src/api.py ./src/
COPY src/predict.py ./src/

USER mluser
EXPOSE 8000
CMD ["uvicorn", "src.api:app", "--host", "0.0.0.0", \
     "--port", "8000"]
\end{minted}
\end{codeblock}

\begin{bestpractice}[title=\textbf{Multi-stage build benefits}]
\begin{itemize}
  \item The final image contains only runtime dependencies---no
    compilers, build tools, or training data.
  \item A typical ML image drops from 5--8\,GB to 500\,MB--1\,GB.
  \item Smaller images mean faster deployments and reduced attack
    surface.
\end{itemize}
\end{bestpractice}

\section{Docker Compose for ML Stacks}
\label{sec:docker-compose}
\index{Docker Compose}

Docker Compose orchestrates multi-container applications with a single
YAML file.

\begin{codeblock}[title=\textbf{docker-compose.yml --- ML application stack}]
\begin{minted}{yaml}
version: "3.9"

services:
  # ML API service
  ml-api:
    build:
      context: .
      dockerfile: Dockerfile
    ports:
      - "8000:8000"
    volumes:
      - model-store:/app/models
    environment:
      - MODEL_PATH=/app/models/model.joblib
      - LOG_LEVEL=info
    depends_on:
      db:
        condition: service_healthy
    deploy:
      resources:
        reservations:
          devices:
            - driver: nvidia
              count: 1
              capabilities: [gpu]

  # PostgreSQL for metadata
  db:
    image: postgres:16-alpine
    environment:
      POSTGRES_DB: mlops
      POSTGRES_USER: mluser
      POSTGRES_PASSWORD: ${DB_PASSWORD}
    volumes:
      - pg-data:/var/lib/postgresql/data
    healthcheck:
      test: ["CMD-SHELL", "pg_isready -U mluser"]
      interval: 5s
      timeout: 5s
      retries: 5

  # MLflow tracking server
  mlflow:
    image: ghcr.io/mlflow/mlflow:v2.10.0
    ports:
      - "5000:5000"
    command: >
      mlflow server
      --backend-store-uri postgresql://mluser:${DB_PASSWORD}@db/mlops
      --default-artifact-root /mlflow/artifacts
      --host 0.0.0.0
    volumes:
      - mlflow-artifacts:/mlflow/artifacts
    depends_on:
      db:
        condition: service_healthy

volumes:
  model-store:
  pg-data:
  mlflow-artifacts:
\end{minted}
\end{codeblock}

\begin{codeblock}[title=\textbf{Docker Compose commands}]
\begin{minted}{bash}
# Start all services
docker compose up -d

# View logs
docker compose logs -f ml-api

# Rebuild after code change
docker compose up -d --build ml-api

# Stop all services
docker compose down

# Stop and remove volumes (destructive)
docker compose down -v
\end{minted}
\end{codeblock}

\section{Best Practices}
\label{sec:docker-best-practices}

\begin{bestpractice}[title=\textbf{Docker best practices for ML}]
\begin{enumerate}
  \item \textbf{Use \file{.dockerignore}}: exclude data, notebooks,
    \file{.git}, \file{\_\_pycache\_\_}, virtual environments.
  \item \textbf{Layer caching}: copy dependency files before source code.
  \item \textbf{Non-root user}: always switch to a non-root user with
    \texttt{USER}.
  \item \textbf{Pin base image versions}: use
    \texttt{python:3.11.7-slim}, not \texttt{python:latest}.
  \item \textbf{Minimize layers}: combine related \texttt{RUN}
    commands with \texttt{\&\&}.
  \item \textbf{Clean up}: remove apt caches with
    \cli{rm -rf /var/lib/apt/lists/*}.
  \item \textbf{Use \texttt{--no-cache-dir}}: prevent pip from storing
    downloaded packages in the image.
  \item \textbf{Health checks}: add \texttt{HEALTHCHECK} for
    production images.
  \item \textbf{Label images}: use \texttt{LABEL} for metadata
    (version, maintainer).
\end{enumerate}
\end{bestpractice}

\begin{codeblock}[title=\textbf{.dockerignore}]
\begin{minted}{text}
.git
.gitignore
__pycache__
*.pyc
.venv/
data/raw/
notebooks/
*.ipynb
.env
.dockerignore
Dockerfile
docker-compose.yml
README.md
\end{minted}
\end{codeblock}

\begin{antipattern}[title=\textbf{Docker anti-patterns for ML}]
\begin{itemize}
  \item Using \texttt{latest} tag for base images---builds become
    non-reproducible.
  \item Running containers as root in production.
  \item Copying the entire project tree including data and notebooks.
  \item Installing dev tools (Jupyter, pytest) in production images.
  \item Storing secrets in the Dockerfile or image layers.
  \item Building images without \file{.dockerignore}---context upload
    takes minutes.
  \item Using a single monolithic container for training, serving, and
    monitoring.
\end{itemize}
\end{antipattern}

\section{Mini-project: Containerized ML Application}
\label{sec:mini-project-ch7}

\begin{miniproject}[title=\textbf{Mini-project 7: Docker for ML}]
\begin{enumerate}
  \item Write a \file{Dockerfile} for the structured ML project from
    Chapter~\ref{ch:introduction}.
  \item Use multi-stage builds: one stage for training, one for
    serving.
  \item Create a \file{.dockerignore} file.
  \item Write a \file{docker-compose.yml} that includes the ML API,
    a PostgreSQL database, and an MLflow tracking server.
  \item Verify that \cli{docker compose up} starts all services and
    that the API responds to prediction requests.
  \item Measure the image size before and after multi-stage
    optimization.
\end{enumerate}
\end{miniproject}

\section{Exercises}
\label{sec:exercises-ch7}

\begin{exercise}[$\bigstar$ --- Basic Dockerfile]
Write a Dockerfile for a scikit-learn project that:
\begin{enumerate}
  \item Uses \texttt{python:3.11-slim} as the base image
  \item Installs dependencies from \file{requirements.txt}
  \item Copies the source code
  \item Runs the training script as the default command
\end{enumerate}
Build the image and verify it runs correctly.
\end{exercise}

\begin{exercise}[$\bigstar\bigstar$ --- Multi-stage optimization]
Starting from a single-stage Dockerfile for a PyTorch project:
\begin{enumerate}
  \item Measure the original image size with \cli{docker images}
  \item Convert to a multi-stage build separating build and runtime
  \item Add a \file{.dockerignore} file
  \item Compare the final image sizes and document the savings
\end{enumerate}
\end{exercise}

\begin{exercise}[$\bigstar\bigstar\bigstar$ --- Full ML stack with Compose]
Create a complete Docker Compose stack for an ML project:
\begin{enumerate}
  \item ML API (FastAPI) with GPU support
  \item PostgreSQL database for predictions logging
  \item MLflow tracking server backed by PostgreSQL
  \item Prometheus for metrics collection
  \item Grafana dashboard for monitoring
\end{enumerate}
All services must have health checks, use named volumes for
persistence, and communicate through a custom Docker network.
\end{exercise}

\section{Cheatsheet}

\begin{cheatsheet}[title=\textbf{Chapter 7 Summary}]
\begin{tabular}{@{}p{4cm}p{8.5cm}@{}}
\toprule
\textbf{Concept} & \textbf{Key Point} \\
\midrule
Image vs Container & Image = template (read-only); Container = running instance \\
\cli{docker build} & Build image from Dockerfile \\
\cli{docker run} & Create and start a container \\
Multi-stage & Separate build/runtime for smaller images \\
Layer caching & Copy \file{requirements.txt} before source code \\
\file{.dockerignore} & Exclude data, notebooks, \file{.git} \\
Non-root user & \texttt{USER mluser} for security \\
Docker Compose & Multi-container orchestration via YAML \\
GPU access & \cli{docker run --gpus all} or Compose \texttt{deploy} \\
\bottomrule
\end{tabular}
\end{cheatsheet}
